\section{Counterparty exposure and valuation adjustments}

A bilaterally cleared derivative can be worth something to either side, so the loss from a counterparty's default depends on what the netting set is worth when the default happens. These entries measure that future exposure, with netting and collateral, price it as a credit valuation adjustment and show how dependence between exposure and default changes the price. The bond case, where exposure is known, is \hyperref[meth:risky-bond-cva]{valuing a risky bond}. Regulatory capital for counterparty risk is a separate calculation, set out in \hyperref[rule:counterparty-credit-cva-capital]{counterparty and CVA capital}. Numbers are checked in \texttt{notes/artefacts/counterparty-xva/counterparty-xva-checks.py}.

\begin{qfdefinition}[Counterparty exposure: expected exposure, potential future exposure and expected positive exposure, with netting, collateral and the margin period of risk]\label{def:counterparty-exposure-measures}
\qffield{Statement} If the counterparty defaults when the netting set has value $V$ to the bank, the bank is an unsecured creditor for $V$ when $V>0$ and owes $-V$ in full when $V<0$. Ignoring the bid--ask cost of replacing the trades, the exposure is $\max(V,0)$, with $V$ the net value of all trades under the master agreement \citep[chap.~18, sec.~18.1, pp.~394--395]{risk-management-financial-institutions-2023}. With collateral $C$ held by the bank (negative if the bank has posted), $E=\max(V-C,0)$ (eq.~(18.4), p.~396). Netting turns trades worth $+10$, $+30$ and $-25$ million into one exposure of 15 instead of 40 \citep[chap.~24, sec.~24.7, p.~574]{s7-hulloptions}.

The exposure measures are taken across simulated values at each future date \citep[chap.~7, sec.~7.2, pp.~114--119]{gregory-xva-challenge-2015}. \emph{Expected future value} is the mean of $V$. \emph{Expected exposure} (EE) is the mean of $\max(V,0)$, above the expected future value as an option is above a forward. \emph{Potential future exposure} (PFE) is a high quantile of exposure, and maximum PFE its largest value over time. \emph{Expected positive exposure} (EPE) is the time-weighted average of EE, and the regulators' effective EPE averages a non-decreasing EE over one year. The negative counterparts (NEE, ENE) are the counterparty's exposure to the bank. For a normal $V$ with mean 2 and standard deviation 2, EE is 2.17 and the 99 percent PFE 6.65; with standard deviation 4 they are 2.79 and 11.31 (p.~116).

\emph{Collateral and the margin period of risk.} A defaulting counterparty stops posting some days before close-out, so the collateral at default reflects the value one margin period of risk (MPR) earlier, typically ten days (RMFI, sec.~18.2.1, pp.~396--398). Hull's four cases with a two-way zero-threshold agreement: values $(50,45)$ at default and twenty days before give exposure 5; $(50,55)$ give 0; $(-50,-45)$ give 0; $(-50,-55)$ give 5, the excess collateral posted and not returned (Example~24.4, pp.~573--574).

\qffield{Notation} $V(t)$ ex-coupon value of the netting set to the bank; $C(t)$ collateral; $\delta$ the MPR; $\mathrm{EE}(t)=E[\max(V(t),0)]$; $\mathrm{PFE}_\alpha(t)$ the $\alpha$ quantile of exposure. Interpretation: under a two-way cash agreement with daily margin and zero threshold, $C(t)=V(t-\delta)$, so $E_c(t)=\max(V(t)-V(t-\delta),0)$, never more than $|V(t)-V(t-\delta)|$.

\qffield{Conditions} Profiles depend on the product. A forward's EE grows like $\sqrt t$; a swap's rises and then falls as payments run off, roughly as $(T-t)\sqrt t$ with a peak near $T/3$ (Gregory, secs.~7.3.2--7.3.3, pp.~120--121). A cross-currency swap keeps rising because the notionals are exchanged at maturity (RMFI, p.~405). Collateral can increase exposure: posted collateral that is not returned is lost (Gregory, p.~133). For a short CDS protection position, posting is so much more likely than receiving that a two-way agreement raises EE (sec.~11.3.6, pp.~258--259). A payment that falls inside the MPR creates a spike, because collateral for it has not arrived (secs.~11.2.5 and~11.3.2, pp.~250--254). For limits, PFE is a scenario measure and in principle needs real-world dynamics, while exposure for pricing is simulated risk-neutrally (RMFI, p.~398).

\qffield{Practice} Checked example (synthetic swaps; Hull--White on the GSW curve of 30 September 2026, $a=0.1$, $\sigma=1$ percent, as in \hyperref[mod:affine-short-rate-models]{affine short-rate models}). An at-market five-year annual payer swap on 100 notional has $V(0)=0$, so $\mathrm{EE}(0)=0$. On 40{,}000 paths EE peaks at 1.82 near two years and is 0.78 a month before maturity. The 97.5 percent PFE at the peak is 7.93 (bootstrap standard error 0.05). Netting with an at-market three-year receiver cuts EE at 1.46 years from 2.15 to 0.68, and the netted exposure never exceeds the sum of the stand-alone exposures on any path. With $\delta=10$ business days and mid-month exposure dates, so that no payment falls inside a margin period, the collateralised EE averages about 0.12 over the dates, against an uncollateralised peak of 1.82.
\begin{itemize}
\item On 52{,}866 path-dates the bank had posted and the value turned positive, so collateral raised exposure there, although mean collateralised EE stayed below EE at every date.
\item A one-way agreement, in which only the bank receives, never raises exposure on any path.
\item The approximation $\sigma_{\Delta V}/\sqrt{2\pi}$, where $\sigma_{\Delta V}$ is the standard deviation of $V(t)-V(t-\delta)$, understates the collateralised EE by 6--8 percent on average at 5, 10 and 20 days. The changes are not normal with mean zero. Gregory writes it as $0.4\,\sigma\sqrt{\text{MPR}}$ (sec.~11.3.3, p.~255).
\end{itemize}

\qffield{Pitfalls} Adding stand-alone exposures across trades under one master agreement. Assuming collateral can only reduce exposure. Choosing exposure dates that straddle payments without modelling the payment during the margin period. Reading EPE as the exposure on a given date.

\qffield{See also} \hyperref[meth:credit-valuation-adjustment]{CVA}; \hyperref[def:wrong-way-risk]{wrong-way risk}; \hyperref[mech:otc-clearing-collateral]{clearing and collateral}; \hyperref[meth:portfolio-var-methods-decomposition]{value at risk}.

\end{qfdefinition}

\begin{qfmethod}[Credit valuation adjustment: unilateral CVA from expected exposure and default probabilities, the independence it needs, DVA, and funding, margin and capital adjustments in outline]\label{meth:credit-valuation-adjustment}
\qffield{Statement} The value of a bilateral portfolio allowing for defaults is $f_{nd}-\mathrm{CVA}+\mathrm{DVA}$, where $f_{nd}$ is the no-default value \citep[chap.~9, sec.~9.1, pp.~216--218]{s7-hulloptions}. Dividing the life into intervals, $\mathrm{CVA}=\sum_i(1-R)\,q_i\,v_i$. Here $q_i$ is the risk-neutral probability that the counterparty defaults in interval $i$, $v_i$ the present value of expected exposure at its midpoint and $R$ the recovery rate. With credit spreads $s_i$, $q_i=e^{-s_{i-1}t_{i-1}/(1-R)}-e^{-s_it_i/(1-R)}$ \citep[chap.~18, sec.~18.2, pp.~395--396, eqs.~(18.1)--(18.3)]{risk-management-financial-institutions-2023}. The formula needs no simulated defaults, only simulated exposures \citep[chap.~14, sec.~14.2.3, p.~311]{gregory-xva-challenge-2015}.

\emph{What the factorization needs} (Interpretation). In general $\mathrm{CVA}=(1-R)\,E^Q[D(0,\tau)E(\tau)\mathbf 1_{\{\tau\le T\}}]$, where $D$ is the bank-account discount factor. Splitting it into $(1-R)\sum_i\mathrm{EE}^*(t_i)\,q_i$ with $\mathrm{EE}^*(t)=E^Q[D(0,t)E(t)]$ needs default to be independent of the \emph{discounted} exposure. A sufficient condition is that the default time is independent of rates and market factors jointly and $R$ is constant. Default correlated with rates breaks the factorization even if it is independent of the swap's value. $\mathrm{EE}^*(t)$ equals $P(0,t)$ times EE under the $t$-forward measure. It equals $P(0,t)$ times EE under $Q$ only when rates are deterministic or the discount factor and exposure are uncorrelated (Gregory, sec.~14.2.6, p.~313).

\emph{DVA and other adjustments.} DVA is the counterparty's CVA, computed from the bank's own default probability and the negative exposure; it rises as the bank's credit worsens (Hull, pp.~217--218; Gregory, sec.~14.6.4, pp.~327--328). FVA prices funding an uncollateralised position: a positive value must be funded, a negative one is a source of funds (Gregory, sec.~15.2.1, pp.~341--342). MVA prices funding initial margin, and KVA the capital a trade requires. Financial economists argue that marginal, not average, funding and capital costs belong in these (Hull, secs.~9.2--9.3, pp.~219--223; RMFI, sec.~18.8, pp.~407--408).

\qffield{Notation} $q_i=\mathrm{PD}(t_{i-1},t_i)$; $\mathrm{EE}^*$ discounted expected exposure; $\tau$ default time. Hull's OFOD defines $v_i$ to include $1-R$, RMFI outside the sum: the same gold forward gives $v_1=92.67$ and $132.38$ in the two books, and both give CVA 5.77 (Example~24.6, pp.~576--577; Example~18.3, pp.~406--407).

\qffield{Conditions} CVA is calculated on the whole netting set; a new trade's incremental CVA comes from rerunning the stored paths with it added (RMFI, sec.~18.3, pp.~399--400). Risk-neutral default probabilities are used for pricing and accounting, although most counterparties have no traded CDS (Gregory, sec.~14.2.7, p.~314). Gregory's rough form $\mathrm{CVA}\approx\mathrm{EPE}\times s$, with $s$ the counterparty's credit spread, separates the credit and market parts. It is a running charge per year, in the units of the spread, with an undiscounted EPE treated as constant; it is not the upfront CVA (sec.~14.2.5, eq.~(14.4) and nn.~4--7, p.~313). Interpretation: the upfront value is a running spread times the risky annuity $A=\sum_i\Delta t_i\,P(0,t_i)\,S(t_i)$, the value of 1 a year paid while the counterparty survives, so $\mathrm{CVA}_{\text{upfront}}\approx\mathrm{EPE}\times s\times A$; Gregory converts the other way, dividing the upfront CVA by $A$.

\qffield{Practice} \emph{Reproduced}: Gregory's square-root exposure ($\mathrm{EE}=1$ percent of $\sqrt t$, 300 basis points, LGD 60 percent) gives 0.200 percent of notional at 20 quarterly dates, 0.193 percent exactly and 0.192 percent with end-point averaging (sec.~14.2.4, p.~312). As a running spread, 0.200 percent divided by his risky annuity of 4.42 (zero rates) is 4.52 basis points a year. The approximation gives $1.54\%\times300=4.62$ basis points a year (p.~313); times 4.42 that is about 0.204 percent upfront, while read as upfront it would be about a quarter of the CVA. The annuity is reproduced by $(1-e^{-hT})/h$ with $h=s/\mathrm{LGD}=5$ percent and $T=5$; Gregory's formula is in his online Appendix 14B, which was not read. RMFI's gold forward gives 5.77, so its value falls from 90.48 to 84.71.

\emph{Checked} (the swaps of \hyperref[def:counterparty-exposure-measures]{exposure measures}, hazard 2 percent independent of rates, $R=40$ percent, monthly intervals): CVA per 100 notional is 0.0601 uncollateralised and 0.0066 with a two-way ten-day agreement. $E^Q[D(0,t)E(t)]$ and $P(0,t)\,\mathrm{EE}(t)$ differ by up to 0.046, about 2.5 percent of the peak exposure. Discounting outside the expectation is therefore not exact when rates are stochastic.

\qffield{Pitfalls} Pulling the discount factor out of the expectation in a rates product. Reading $\mathrm{EPE}\times s$ as an upfront CVA rather than a charge per year. Using historical default probabilities in a price. Calculating CVA trade by trade under netting. Counting a DVA gain as a source of capital, which regulators exclude (RMFI, p.~403).

\qffield{See also} \hyperref[def:counterparty-exposure-measures]{exposure measures}; \hyperref[def:wrong-way-risk]{wrong-way risk}; \hyperref[meth:risky-bond-cva]{risky bond CVA}; \hyperref[meth:cds-pricing]{CDS pricing}; \hyperref[rule:counterparty-credit-cva-capital]{regulatory capital}.

\end{qfmethod}

\begin{qfdefinition}[Wrong-way and right-way risk: dependence between exposure and the counterparty's default, its sources and its effect on CVA]\label{def:wrong-way-risk}
\qffield{Statement} Wrong-way risk is positive dependence between the counterparty's default probability and the bank's exposure: default tends to come when exposure is high. Right-way risk is the opposite \citep[chap.~18, sec.~18.5, pp.~401--402]{risk-management-financial-institutions-2023}. With dependence, CVA must use the exposure conditional on default in each interval, in place of the unconditional one \citep[chap.~17, sec.~17.3.1, eq.~(17.1), p.~380]{gregory-xva-challenge-2015}. Here $E(t)$ is undiscounted, so the discount factor goes inside the conditional expectation:
\[ \mathrm{CVA}=(1-R)\sum_iE^Q\big[D(0,t_i)E(t_i)\,\big|\,t_{i-1}<\tau\le t_i\big]\,q_i,\qquad q_i=Q(t_{i-1}<\tau\le t_i), \]
the date-grid approximation of $(1-R)\,E^Q[D(0,\tau)E(\tau)\mathbf 1_{\{\tau\le T\}}]$. Gregory's EE is already discounted (sec.~14.2.3, p.~311), and he conditions on $t_i$ being the default time, which on the grid is default in the interval ending at $t_i$. With deterministic rates $D(0,t_i)=P(0,t_i)$ comes outside; under independence the conditional expectation is the unconditional $\mathrm{EE}^*$ of \hyperref[meth:credit-valuation-adjustment]{CVA}.

\emph{Sources.} Buying protection from a counterparty whose credit is correlated with the reference entity, or a put on a bank's shares from another bank. Paying a sovereign's currency forward, or a corporate paying fixed in a swap when rate cuts come with recessions. A company speculating with many similar trades. Hedges are usually right-way: a mining company that hedges owes money when the metal price is high and its business is strong (Gregory, sec.~17.2.2, pp.~378--379; RMFI, p.~402). \emph{General} wrong-way risk comes from macroeconomic links and can be modelled. \emph{Specific} wrong-way risk comes from structural links, is hard to detect in data and is better avoided (Gregory, sec.~17.2.3, p.~379).

\qffield{Notation} Latent-variable form (Interpretation, the structural approach of Gregory, sec.~17.4.2, p.~385): default by the date if $Y<\Phi^{-1}(\mathrm{PD})$ with $Y$ standard normal, exposure from $V$, and $\rho=\mathrm{Corr}(V,Y)$. Lower-tail default with $\rho<0$ pairs default with high $V$, so it is wrong-way. An upper-tail convention works equally if the sign is changed throughout.

\qffield{Conditions} Hazard-rate models that correlate credit spreads with market factors produce only weak dependence even at correlations near $-90$ percent. Structural mappings of default time and exposure can produce strong effects (Gregory, sec.~17.4, pp.~383--385). Wrong-way risk matters more for a better-rated counterparty, because its default is the bigger surprise (sec.~17.3.2, pp.~381--382). A zero measured correlation does not prove independence (p.~380, n.~3).

\qffield{Practice} Checked example (synthetic): one exposure date, $V\sim N(0,1)$, PD 5 percent, $R=40$ percent, unit discounting. Bivariate-normal quadrature gives CVA 0.03253, 0.02097, 0.01197, 0.00550 and 0.00159 at $\rho=-0.5$, $-0.25$, 0, 0.25 and 0.5, each within three standard errors of 200{,}000 simulated draws. The independent value is $(1-R)\,\mathrm{PD}/\sqrt{2\pi}=0.01197$. The strongest wrong-way case is 2.7 times it, the strongest right-way case 13 percent of it. Discounting matters once $D\ne1$: a one-year exposure of 100 with PD 5 percent, $R=40$ percent and a deterministic continuously compounded rate of 5 percent has $\mathrm{CVA}=0.6\times0.05\times100\,e^{-0.05}=2.8537$, not 3.

\qffield{Pitfalls} Reading a correlation sign without fixing which tail defaults. Treating historical correlations as evidence of no wrong-way risk when the link is causal or specific. Posting collateral in a currency or security that falls when the exposure rises (Gregory, sec.~17.3.3, pp.~382--383).

\qffield{See also} \hyperref[meth:credit-valuation-adjustment]{CVA}; \hyperref[def:counterparty-exposure-measures]{exposure measures}; \hyperref[meth:cds-strategies]{CDS strategies}; \hyperref[rule:counterparty-credit-cva-capital]{regulatory capital}.

\end{qfdefinition}

\begin{qfrule}[Regulatory capital for derivatives counterparty risk: SA-CCR exposure at default, default-risk capital and the separate CVA-risk capital, with their EU, UK and US implementation]\label{rule:counterparty-credit-cva-capital}
\qffield{Statement} Bank capital for a derivatives counterparty involves three different quantities.
\begin{itemize}
\item \emph{Exposure measurement.} A bank without internal-model approval measures each netting set's exposure at default with the standardised approach (SA-CCR): $\mathrm{EAD}=\alpha(\mathrm{RC}+\mathrm{PFE})$ with $\alpha=1.4$ \citep[CRE52.1]{basel-framework-ccr-mar-2026}. Unmargined, the replacement cost is $\mathrm{RC}=\max(V-C,0)$ (CRE52.10), and $\mathrm{PFE}=\text{multiplier}\times\mathrm{AddOn}$ (CRE52.20).
\item \emph{Default-risk capital.} The EAD then enters the standardised or IRB credit-risk approach as the exposure and is risk-weighted for the counterparty's default (CRE51.7 and~51.13).
\item \emph{CVA-risk capital.} A separate charge covers losses from changes in CVA as counterparty spreads and market factors move, with RWA $=12.5\times$ the requirement (MAR50.1 and~50.4). Regulatory CVA excludes the bank's own default (MAR50.3).
\end{itemize}

\emph{SA-CCR mechanics.} For an interest-rate trade the adjusted notional is the notional times the supervisory duration $\mathrm{SD}=(e^{-0.05S}-e^{-0.05E})/0.05$, with $S$ and $E$ the start and end of the referenced period (CRE52.34). For an unmargined trade it is multiplied by the maturity factor $\sqrt{\min(M,1)}$, with $M$ the remaining maturity in years floored at ten business days, $10/250$ years under a 250-day convention (CRE52.48--52.49), and by a supervisory delta of $\pm1$ for a linear trade (CRE52.39). The floor matters for short trades: with five business days left the factor is 0.2, not 0.141. Trades are netted by currency in three maturity buckets, with partial offset between buckets, and the hedging set's add-on is 0.5 percent of its effective notional (CRE52.57). Excess collateral or a negative value reduces PFE but not below 5 percent of the add-on:
\[ \text{multiplier}=\min\Big\{1;\;0.05+0.95\exp\Big(\frac{V-C}{2\times0.95\times\mathrm{AddOn}}\Big)\Big\}\quad\text{(CRE52.23).} \]

\emph{The basic approach to CVA risk} (BA-CVA, reduced version, without hedges). A counterparty's stand-alone requirement is
\[ \mathrm{SCVA}_c=\frac1\alpha\,\mathrm{RW}_c\sum_{NS}M_{NS}\,\mathrm{EAD}_{NS}\,\mathrm{DF}_{NS},\qquad \mathrm{DF}=\frac{1-e^{-0.05M}}{0.05M}, \]
with $\mathrm{RW}_c$ from 0.5 to 12 percent by sector and credit quality (MAR50.15--50.16). Across counterparties
\[ K_{\text{reduced}}=\sqrt{\Big(\rho\sum_c\mathrm{SCVA}_c\Big)^2+(1-\rho^2)\sum_c\mathrm{SCVA}_c^2},\qquad\rho=50\ \text{percent}, \]
and the requirement is $0.65\,K_{\text{reduced}}$ (MAR50.14). The standardised approach (SA-CVA), based on CVA sensitivities, needs supervisory approval (MAR50.7). A bank with non-centrally cleared derivatives of no more than EUR~100 billion notional may instead set CVA capital equal to its counterparty-credit-risk capital (MAR50.9).

\qffield{Notation} $V$ value of the netting set to the bank; $C$ haircut collateral held; $M$ remaining (MF) or effective (BA-CVA) maturity in years; $\alpha$ converts effective expected positive exposure (EEPE) to EAD, which is why BA-CVA divides by it (MAR50.15, n.~4). $\mathrm{SD}$ is a continuously discounted annuity at 5 percent (Interpretation).

The product $M\,\mathrm{EAD}\,\mathrm{DF}/\alpha=M\cdot\mathrm{DF}\cdot\mathrm{EEPE}$ is a supervisory proxy for the \emph{area} under the discounted expected-exposure profile (MAR50.15, n.~3), in currency times years. It is not the discounted expected exposure $\mathrm{EE}^*(t)$ of \hyperref[meth:credit-valuation-adjustment]{CVA}, which is a value in currency at one date. Interpretation: the proxy stands in for $\int_0^M\mathrm{EE}^*(t)\,dt$, with a flat EEPE in place of the modelled profile and a fixed 5 percent rate in place of the discount curve. Multiplying by the risk weight, a supervisory spread move per year, gives SCVA in currency.

\qffield{Conditions} Basel dates: SA-CCR became effective on 1 January 2017 \citep[sec.~D, p.~3]{bcbs-sa-ccr-2014}; the current SA-CCR and CVA chapters are the versions effective 1 January 2023 \citep[CRE52 and MAR50]{basel-framework-ccr-mar-2026}. Implementation, as of 21 October 2026:
\begin{itemize}
\item \emph{EU.} SA-CCR applies from 28 June 2021 \citep[Art.~3(2)]{eu-2019-876-crr2}. CRR3 replaced the CVA approaches with standardised, basic and simplified ones from 1 January 2025 \citep[Art.~1(199) and Art.~2]{eu-2024-1623-crr3}, but kept the exemptions for, among others, non-financial counterparties below the EMIR clearing threshold, intragroup trades and sovereign-type counterparties \citep[Art.~382(4)]{eu-2013-575-crr}.
\item \emph{UK.} SA-CCR applies from 1 January 2022 \citep[p.~8]{pra-ps17-21-basel-2021}. The Basel 3.1 CVA approaches (an alternative approach for small users, BA-CVA and SA-CVA) take effect on 1 January 2027 and remove existing exemptions, with transitional arrangements \citep[paras~1.8 and~1.13]{pra-ps17-23-basel31-2023}; \citep[para.~1.23]{pra-ps1-26-basel31-2026}.
\item \emph{US.} SA-CCR took effect on 1 April 2020, mandatory for advanced-approaches banks from 1 January 2022 \citep[p.~4362]{us-sa-ccr-rule-2020}. Alpha is 1 for a commercial end-user \citep[12 CFR 217.132(c)(5)(iv)]{ecfr-12cfr217-2026}. CVA capital is still the simple or advanced CVA approach of the advanced approaches (217.132(e)); BA-CVA and SA-CVA are only proposed \citep[p.~14955]{us-capital-proposal-2026}.
\end{itemize}

\qffield{Practice} Checked example (\texttt{counterparty-capital-check.py}; original inputs). An unmargined four-year payer swap on 100 notional is worth 2 to the bank, with no collateral.
\begin{itemize}
\item \emph{SA-CCR}: SD 3.6254, adjusted notional 362.54, add-on $0.005\times362.54=1.8127$, RC 2 and multiplier 1, so $\mathrm{EAD}=1.4\times(2+1.8127)=5.3378$. With 4 of cash collateral held, RC is 0, the multiplier 0.5815 and EAD 1.4758. For a US commercial end-user ($\alpha=1$) the uncollateralised EAD is 3.8127.
\item \emph{Default-risk capital} at an illustrative 100 percent risk weight: RWA 5.3378, capital 0.4270 at 8 percent.
\item \emph{CVA-risk capital} for an investment-grade industrial ($\mathrm{RW}=3$ percent), taking $M=4$ for illustration: DF 0.9063, SCVA 0.4147, requirement $0.65\times0.4147=0.2695$, RWA 3.3692. The exposure area $M\,\mathrm{EAD}\,\mathrm{DF}/\alpha$ is $4\times5.3378\times0.9063/1.4=13.82$ currency-years, an EEPE of 3.81 held for four years at a 5 percent discount; SCVA is $0.03\times13.82$. Two such counterparties need $\sqrt{2.5}=1.58$ times one, not twice.
\end{itemize}

\qffield{Pitfalls} Treating EAD as a capital charge; it is the exposure fed to the credit-risk approach. Adding default-risk and CVA-risk capital as if one measured the other. Reading regulatory CVA as accounting CVA. Quoting the Basel effective date as the date a jurisdiction applies the rule.

\qffield{See also} \hyperref[def:counterparty-exposure-measures]{exposure measures}; \hyperref[meth:credit-valuation-adjustment]{CVA}; \hyperref[def:wrong-way-risk]{wrong-way risk}; \hyperref[mech:otc-clearing-collateral]{clearing and collateral}.

\end{qfrule}
